\section{Experimental design}
\label{sec:experimental-design}

\subsection{Predictions and when they were set}
\label{sec:design-hypotheses}

Each experiment states its predictions in the header of its script, and \cref{tab:hypotheses} lists all \hypTotal{} with the value observed and the outcome, computed from the experiment output by the predicate written beside the statement.
The predictions were written by the author before the first run of the script that tests them, on 7 November 2026.
The repository records a script and its output in one change, so the order of writing and running rests on that statement and not on the history of the repository.
The history does fix the age of the committed documents that the scripts read: the C06 and C07 documents were first baked on 12 October 2026, with the tests of their identities, of the separate overlay lines and of the closing of the bridge, and were re-baked on the engine 1.1.0 on 30 October 2026.
The acceptance criteria that the product's design document states for this subject, that the identities close within the tolerance, that the weighted ECL is at least the average-scenario ECL in the convex regime, that the reference example gives 2.68\% against 1.80\%, and that overlays are separate lines, entered that document on 28 October 2026, and the claim text on 20 and 21 October, after the bakes they describe.
They cannot fail on the committed documents by construction of the bake and are therefore tabulated as reproductions (the rows marked after the bake), not as predictions; one of them, that the weighted ECL exceeds the average-scenario ECL on every committed date, is stated in the product as not claimed, and is tabulated because the committed documents contradict it.
A second kind of reproduction, in the rows that say so, is a prediction written before the run of a script for a result that the product's notebook had already measured with another seed or on another defect class.
The remaining predictions concern results that no committed document holds: the comparators, the rounding, the placement of overlays, the scenario sets, the ten seeds and the structure of the identities.

\subsection{Data and provenance}
\label{sec:design-data}

The first data set is the SBA 7(a) loan book of the U.S. Small Business Administration, a U.S. Government Work released under the FOIA data programme \citep{sba7a}, read by the product's C06 case at six reporting dates between 2006 and 2023.
The outcomes are real; the book is a public twin of a bank loan book, not a bank's own, and it holds small-business loans, not mortgages.
Each date gives a staged book with its allowance and gross carrying amount roll-forwards, its waterfall and its stage summary, produced by the engine on a discrete-time hazard model fitted on data up to the date, with the scenario factors derived from the 2026 supervisory scenario tables of the Federal Reserve.
\cref{tab:data} gives, for each date, the gross total $\gross$, the limit of the test, the Stage 2 share and the largest residual on the committed tables.

The second data set is the synthetic bank of the C07 case, Banco Andino: accounts drawn from a seeded generator and calibrated to the aggregates published by the CMF (portfolio mix, provision index and 90-day delinquency ratio by portfolio), evaluated in six scenarios (base, recession, house price fall, consumer arrears shock, UF inflation spike, combined severe) by the engine's CMF B-1, IFRS 9 and Annex 9 modules.
It is not a real bank and says so in every document it produces.
The committed documents are the pair seed 42 at the full size of \efiveAccounts{} accounts.

The third source is resampled: from each committed C06 table 200 clean books are drawn by multiplying every opening and movement by a factor $1+0.25\,\xi$ with $\xi$ standard normal truncated at $\pm3$, setting each closing to the exact sum, each waterfall line to the sum over stages and each stage ECL to its stage closing.
A resampled book therefore has the shape and size of one of the six dates and closes exactly.
The sweeps of Sections~\ref{sec:method-nonlinearity} and~\ref{sec:theory-convexity} draw their inputs from stated distributions and use no data set.

\subsection{Methods compared and the data they share}
\label{sec:design-comparators}

Every comparison is run by the product on the same cases.
For the detection of planted defects the compared methods are the twelve identities, the stage arithmetic (the six roll-forward identities alone) and the pooled row; each defect, at a given size and location, is judged by all three, and the leave-one-family-out variants of the twelve.
For the false alarms the compared methods are the relative tolerance $10^{-12},10^{-9},10^{-6},10^{-4}$ of $\gross$ and the absolute tolerance $10^{-6},10^{-2},1$ units, judged on the same six committed tables and the same 1,200 resampled books, and on the same books after rounding.
For the overlays the compared bookings are the correct one and five wrong ones, judged by the three methods on the same six dates and three shares of the allowance, 2\%, 10\% and 25\%.
For the non-linearity the compared methods are five scenario sets judged against the same quadrature on the same \efourDesignDraws{} draws, and the index is computed by the engine's own test row.
For the regimes the compared methods are the three regimes themselves, run on the same generated book, in each of the six scenarios and ten seeds.
No result of a published method is taken from a paper: each one is recomputed by the engine or the product's scripts.

\subsection{Metrics}
\label{sec:design-metrics}

The detection rate of a method at a size is the share of cases flagged, with the Wilson interval \citep{wilson1927,brown2001}
\begin{equation}
\Big[\frac{\hat p+z^{2}/2n\pm z\sqrt{\hat p(1-\hat p)/n+z^{2}/4n^{2}}}{1+z^{2}/n}\Big],\qquad z=1.96.
\label{eq:wilson}
\end{equation}
The detection limit is the ratio $\delta^{*}_j/\tlim$ of the bisected size to the limit.
The false-alarm rate is the share of clean books flagged.
The non-linearity is measured by the index $\nu$ of \eqref{eq:index}, by the ratio of the second-order formula \eqref{eq:second-order} to $\nu$, and for a scenario set by the relative error $\varepsilon_d=(P_d-P^{*})/P^{*}$, summarised by its quantiles and by the median of $|\varepsilon_d|$ with a bootstrap interval \citep{efron1979}.
The regimes are compared by the ratios $T_I/T_B$ and $T_A/T_B$, summarised over seeds by their mean and the $t$ interval of the mean with nine degrees of freedom, and by the largest bridge residual.
Two further measures describe the structure of the identities and are not random: the rank of the matrix and the number of entries each equation sees.
The agreement of the engine with the signed matrix is the largest absolute deviation of an engine residual from the prediction of the matrix, divided by the size of the perturbation, over all entries and dates.
A prediction is held when its predicate on these measures is true on the experiment output, and refuted otherwise; no measure is rounded before the predicate is applied.

\subsection{Repetitions, seeds and statistics}
\label{sec:design-statistics}

The design follows the account of a simulation study by aims, data-generating mechanisms, methods, estimands and performance measures \citep{morris2019}.
The first experiment plants defects at \eoneReps{} random locations in each of six dates and eight classes, \eoneCases{} cases per class, at nine sizes from 0.25 to 100 times the limit, with the same location at every size, and bisects \eoneBisected{} cases per class.
The second draws \etwoReps{} resampled books per date.
The third draws \ethreeDraws{} overlay-set pairs and enumerates the 18 placement cases per booking.
The fourth draws \efourDraws{} factor sets for the sign of the index and \efourDesignDraws{} forecasts for the scenario sets, with 2,000 bootstrap resamples of each design.
The fifth runs the generator of the synthetic bank for \efiveSeeds{} seeds (42 and 1 to 9) at the committed size and for seed 42 at six sizes.
The per-run seed of each draw is derived from a base seed and the indices of the experiment, the case and the draw, and is stated in the script, so that a rerun gives byte-identical outputs.
Run times are those the scripts measured: the first experiment \eoneSeconds{} seconds, the fourth \efourSeconds{} seconds, and one build of the six scenarios of the synthetic bank between \efiveSecLow{} and \efiveSecHigh{} seconds at the committed size.
